\section{Methodology and Results}
\label{sec:methodology}

\subsection{Data Description and Pre-Processing}
\label{subsec:data}

The dataset used for training was obtained from historical traded option contracts on the \textit{SPY ETF}\footnote{An exchanged traded fund replicating the S\&P500 index} from January 2021 to December 2022, with a total of 68 unique days.
The dataset contains various columns on traded options, indexed by date, including strike prices, future contract prices, maturities, implied volatilities, among others.
Relevant to our study, the principal columns used were:
\begin{itemize}
    \item \textbf{Underlying}: Price of the underlying asset.
    \item \textbf{Strike}: Option's strike price.
    \item \textbf{Expiration}: Days remaining until the option's expiration.
    \item \textbf{Implied Volatility}: Value for the implied volatility.
    \item \textbf{Date}: Traded date for the contract, used as the preprocessed dataset index.
\end{itemize}

Additionally, two columns were added to be used for the Monte Carlo estimation:
\begin{itemize}
    \item \textbf{Maturity}: Calculated by dividing \textit{Expiration} by 252 (working days in a year) to normalize the maturity time.
    \item \textbf{Risk-free rate}: Risk-free rate obtained from the Federal Reserve Economic Data (FRED) for the corresponding period.
    \item \textbf{Dividend}: Dividend yield paid quarterly by the SPY\@ index.
\end{itemize}

\subsection{Data Preprocessing}
\label{subsec:preprocessing}

The data preprocessing involved several essential steps, predominantly to allow the use of the surface data directly
in the transformer architecture and subsequent generation of the surface for historical data.
Each of the trained transformers (one for each parameter $\alpha, \rho$, or $\nu$) takes four dimensions as input, which are:
\begin{enumerate}
    \item Time to maturity;
    \item The daily parameters $\alpha, \rho$, or $\nu$ calculated using the \textit{BFGS} optimization algorithm;
    \item $\mathbf{1}$ if the parameter was calculated for the current day, $\mathbf{0}$ if it was calculated for a fixed set of maturities
    (top 7 most frequent maturities) using the previous day fitted parameters;
    \item Parameter value calculated using the previous day's $\mathbf{P}$, $\mathbf{Q}$ and $\mathbf{R}$ values for same strike and maturity.
\end{enumerate}

The target scores were generated before training using the Monte Carlo algorithm, detailed in \hyperref[subsubsec:monte_carlo]{Section~\ref*{subsubsec:monte_carlo}}.
with $p = 80\%$ and $k = 1 \times 10^{2}$.
For training and testing we also observed that the choice of $\beta = 0.5$ performed best at minimizing the surface's MAE.

\subsection{Results and Surface Comparison}
\label{subsec:results}

A baseline linear interpolation for the daily observed points was also used to compare our trained
model's surfaces by interpolating the points within a fixed grid of strikes and maturity.
No extrapolation was performed due to the nature of the linear model.

The training dataset had a total of $44817$ observations totaling 68 unique days, which lead to the decision of using a leave-one-out cross-validation approach.
The Adam algorithm was used as optimizer, with a max epoch range of $10^{3}$ and an early stopping sensitivity of $10^{-3}$.
The Anderson acceleration was implemented with a maximum iteration limit of 10, a
convergence threshold of $\epsilon=10^{-4}$, damping factor of $\lambda=0.1$, and memory size of $m=5$.
A batch size of 8 was used to avoid memory issues when training the implicit layer.
The network architecture is a transformer encoder block with 8 sequential attention heads,
2 linear layers with 256 hidden units and 2 implicit layers with 256 and 128 hidden units, respectively, and a dropout rate of 0.1 was used for the dropout layers.
An illustration of the architecture is shown in \hyperref[fig:figure1]{Figure~\ref*{fig:figure1}}.

\begin{figure}[h]
    \begin{center}
        \includegraphics[width=0.5\textwidth]{images/losses}
        \caption{Historical loss curve on training session data. Curves were cut for better visualization,
            with the largest episode being 2950 epochs for the $SST_\alpha$ (alpha) network.}
        \label{fig:figure3}
    \end{center}
\end{figure}

\begin{table}[!ht]
    \centering
    \renewcommand{\arraystretch}{1.5} % Adjust the vertical padding
    \resizebox{\columnwidth}{!}{%
        \begin{tabular}{|l|l|l|l|}
            \hline
            Model          & Test loss (RMSE) & Train time & Preprocess time (Monte Carlo Estimation) \\ \hline
            $SST_{\alpha}$ & 1.21             & 4m21s      & 7m50s                                    \\ \hline
            $SST_{\rho}$   & 0.72             & 2m8s       & 6m34s                                    \\ \hline
            $SST_{\nu}$    & 0.75             & 1m43s      & 6m41s                                    \\ \hline
            Total          & -                & 8m12s      & 21m5s                                    \\ \hline
        \end{tabular}%
    }
    \captionsetup{skip=10pt} % Adjust the caption margin
    \caption{Model training metrics for each of the three SST models.}
    \label{tab:training}
\end{table}

For benchmarking, the execution time was measured using the Python \texttt{perf\_counter} method and a heat-up of the model was performed beforehand on the entire dataset.
Separate tests using a linear interpolation with the original points and with the implicit layer were performed to compare both generated surfaces,
shown in \hyperref[sec:surface_comparison]{Appendix~\ref*{sec:surface_comparison}}.
The linear interpolation obtained a MAE of $0.0$ to the observed surface, as expected,
and the proposed model had a MAE of $7 \times 10^{-4}$ for the points used in the Monte Carlo estimation.
While still near the threshold hyperparameter of $1 \times 10^{-4}$, the threshold could be ensured by using a larger maximum iteration limit for the Anderson acceleration,
trading off for a larger training and execution time.

The mean train and test losses using LOO cross-validation are shown in \hyperref[tab:training]{Table~\ref*{tab:training}},
while the execution times shown below in \hyperref[tab:performance]{Table~\ref*{tab:performance}} are on the entire dataset.
The RMSE values for the test data show that the model was able to generalize well to the test data, with the lowest RMSE value for the $\rho$ parameter.

\begin{table}[!ht]
    \centering
    \begin{tabular}{|l|l|l|l|}
        \hline
        Metric                           & Mean time & Standard Deviation \\ \hline
        Execution time on entire dataset & 229.1ms   & 41.2ms             \\ \hline
        Execution time per point         & 2ms       & 1ms                \\ \hline
    \end{tabular}%
    \captionsetup{skip=10pt} % Adjust the caption margin
    \caption{Model performance}
    \label{tab:performance}
\end{table}

